% ============================================================
% Section 124: 硅价带宽度的自由电子模型估算
% ============================================================
\section{固体物理：硅价带宽度的自由电子模型估算}
\label{sec:si-valence-width}

\begin{modelbox}[题目：自由电子模型估算硅的价带宽度]
若把半导体硅材料价带中的电子看作是自由电子，试计算其价带宽度等于多少电子伏特？（硅晶体晶格常数 $a=5.4\,\text{\AA}$。）
\end{modelbox}

\begin{tipbox}[考点快速分析]
\begin{itemize}
    \item \textbf{自由电子模型}：$E(k)=\hbar^2k^2/2m$，价带宽度对应费米能 $E_F$
    \item \textbf{硅晶体结构}：金刚石结构，每个惯用晶胞含 8 个原子，每个原子 4 个价电子
    \item \textbf{费米波矢}：$k_F=(3\pi^2n)^{1/3}$，$n$ 为价电子数密度
    \item \textbf{数值换算}：注意 $\hbar$ 与 $h$ 的使用，以及 eV 与 J 的换算
\end{itemize}
\end{tipbox}

\begin{derivationbox}[标准解题过程]
\textbf{1. 自由电子模型与费米波矢}

将硅的价带电子视为自由电子气，能量色散关系为
\[
E(k)=\frac{\hbar^2k^2}{2m}.
\]
在绝对零度下，电子从 $E=0$ 填充至费米能级 $E_F$，在 $k$ 空间形成半径为 $k_F$ 的费米球。价带宽度即对应于费米能 $E_F=E(k_F)$。

电子数密度 $n=N/V$ 与费米波矢的关系（计入自旋简并）为
\[
n=\frac{k_F^3}{3\pi^2}
\implies k_F=(3\pi^2n)^{1/3}.
\]

\textbf{2. 硅的电子数密度}

硅为金刚石结构，每个惯用晶胞边长为 $a$，内含 8 个硅原子，每个硅原子有 4 个价电子。因此价电子数密度为
\[
n=\frac{8\times4}{a^3}=\frac{32}{a^3}.
\]
代入 $a=5.4\times10^{-10}\,\text{m}$：
\[
n=\frac{32}{(5.4\times10^{-10})^3}\approx\frac{32}{1.575\times10^{-28}}\approx2.03\times10^{29}\,\text{m}^{-3}.
\]

\textbf{3. 费米能与价带宽度}

费米波矢：
\[
k_F=\left(3\pi^2\cdot\frac{32}{a^3}\right)^{1/3}=\frac{(96\pi^2)^{1/3}}{a}\approx\frac{9.82}{5.4\times10^{-10}}\approx1.82\times10^{10}\,\text{m}^{-1}.
\]

费米能（价带宽度）：
\[
E_F=\frac{\hbar^2k_F^2}{2m}=\frac{(1.055\times10^{-34})^2\times(1.82\times10^{10})^2}{2\times9.11\times10^{-31}}\approx2.02\times10^{-18}\,\text{J}.
\]
换算为电子伏特：
\begin{equation}
\boxed{E_F=\frac{2.02\times10^{-18}}{1.602\times10^{-19}}\approx12.6\,\text{eV}.}
\end{equation}

\textbf{4. 讨论}

实际硅的价带宽度由能带理论计算和实验（如 X 射线光电子能谱）测定约为 $12\sim13\,\text{eV}$。自由电子模型的估算值 $12.6\,\text{eV}$ 与实际值惊人地接近，说明从总态密度和能量填充的角度看，自由电子气模型给出了价带整体宽度的良好数量级估计。

然而，自由电子模型无法描述能带结构的细节，如禁带、轻重空穴带的分裂、有效质量的各向异性等。这些需要更精细的能带计算方法（如赝势法、紧束缚法）才能得出。
\end{derivationbox}

\begin{formulabox}[答案汇总]
\begin{center}
\begin{tabular}{ll}
\toprule
\textbf{物理量} & \textbf{答案} \\
\midrule
价电子数密度 $n$ & $\approx2.03\times10^{29}\,\text{m}^{-3}$ \\
费米波矢 $k_F$ & $\approx1.82\times10^{10}\,\text{m}^{-1}$ \\
价带宽度 $E_F$ & $\approx12.6\,\text{eV}$ \\
\bottomrule
\end{tabular}
\end{center}
\end{formulabox}

\begin{insightbox}[物理图像]
\begin{itemize}
    \item 自由电子模型通过费米能估算满带宽度，仅依赖于电子密度这一宏观参数。硅的价带宽度约 $12.6\,\text{eV}$，与实验值一致，说明价电子在整体填充行为上近似自由。
    \item 该简算体现了自由电子模型在估算某些宏观能带参数时的有效性，但其局限性在于无法给出能带结构的细节（如有效质量、各向异性、带隙等）。
\end{itemize}
\end{insightbox}
